\documentclass[aspectratio=169]{beamer}
\usetheme{Madrid}
\usefonttheme{structurebold}
\usefonttheme[onlymath]{serif}

\definecolor{BKTemplate}{rgb}{0.11765, 0.30588, 0.47451}
\usecolortheme[named=BKTemplate]{structure}
\setbeamercolor{title}{fg=BKTemplate,bg=}

\usepackage[utf8]{inputenc}
\usepackage[T1]{fontenc}    % use 8-bit T1 fonts
\usepackage[vietnamese=nohyphenation]{hyphsubst} % fix warning when using vietnamese babel
\usepackage[english, vietnamese]{babel}
\usepackage{url}            % simple URL typesetting
\usepackage{booktabs}       % professional-quality tables
\usepackage{amsfonts}       % blackboard math symbols
\usepackage{microtype}      % microtypography
\usepackage{mathtools}
\usepackage{amsmath}
\usepackage{xcolor}
\usepackage{graphicx}
\usepackage{ifthen}
\usepackage{tikz}
\usetikzlibrary{arrows.meta, positioning, calc}
\usepackage{pifont}

%------------------------------------------------------------
% Thông tin trang tiêu đề
\title[Tên ngắn] % (tùy chọn) hiện ở footer
{\vspace*{1cm}\\\textbf{TÊN ĐỀ TÀI}\\\textbf{Tên đầy đủ của đề tài}\\
\small{Phụ đề / tên tiếng Việt}}

\subtitle{[Hội nghị / buổi báo cáo] $\cdot$ [Lab]}

\author[Họ tên] % (tùy chọn) hiện ở footer
{GVHD: [Học hàm, học vị, họ tên]\\
Học viên: [Họ tên] - [MSHV]}

\institute[HCMUT] % (tùy chọn)
{
  Faculty of Computer Science and Engineering\\
  Ho Chi Minh City University of Technology\\
  Vietnam National University Ho Chi Minh City
}

\date[\the\year] % (tùy chọn)
%------------------------------------------------------------

%------------------------------------------------------------
% Mục lục ở đầu mỗi section, tô đậm section hiện tại.
% Xóa khối này nếu không muốn slide mục lục lặp lại.
\AtBeginSection[]
{
  \begin{frame}
    \frametitle{Nội dung}
    \tableofcontents[currentsection]
  \end{frame}
}
%------------------------------------------------------------

\begin{document}
% Footer chỉ đếm tới hết phần chính (nhãn endmain), không tính slide dự phòng.
\setbeamertemplate{page number in head/foot}{\insertframenumber{} / \pageref{endmain}}

% Nền trang tiêu đề
\setbeamertemplate{background canvas}{
	\includegraphics[width=\paperwidth,height=\paperheight]{images/FrontBackground.pdf}
}
\begin{frame}[plain]
\titlepage
\end{frame}

% Nền cho các trang còn lại
\setbeamertemplate{background canvas}{
	\includegraphics[width=\paperwidth,height=\paperheight]{images/Background.pdf}
}

\begin{frame}
\frametitle{Nội dung}
\tableofcontents
\end{frame}

%---------------------------------------------------------
\section{Giới thiệu}
%---------------------------------------------------------

\begin{frame}
\frametitle{Slide văn bản}
\begin{itemize}
  \item Ý thứ nhất
  \item Ý thứ hai, có \textbf{nhấn mạnh} và công thức $f(x) = x^2$
  \begin{itemize}
    \item Ý con
  \end{itemize}
\end{itemize}

\begin{block}{Tiêu đề khối}
Nội dung khối.
\end{block}
\end{frame}

\begin{frame}
\frametitle{Hai cột}
\begin{columns}[T]
\begin{column}{0.48\textwidth}
\textbf{Cột trái}
\begin{itemize}
  \item Nội dung
\end{itemize}
\end{column}
\begin{column}{0.48\textwidth}
\centering
\includegraphics[width=0.5\textwidth]{images/BK.png}
\end{column}
\end{columns}
\end{frame}

%---------------------------------------------------------
\section{Phương pháp}
%---------------------------------------------------------

\begin{frame}
\frametitle{Bảng}
\begin{center}
\begin{tabular}{@{}lcc@{}}
  \toprule
  Phương pháp & Độ đo A & Độ đo B \\
  \midrule
  Baseline   & 0{,}00 & 0{,}00 \\
  Đề xuất    & \textbf{0{,}00} & \textbf{0{,}00} \\
  \bottomrule
\end{tabular}
\end{center}
\end{frame}

\begin{frame}
\frametitle{Hình vẽ TikZ}
\begin{center}
\begin{tikzpicture}[font=\small, node distance=12mm,
  box/.style={draw=BKTemplate, fill=BKTemplate!8, rounded corners=2pt, inner sep=5pt, align=center}]
  \node[box] (a) {Đầu vào};
  \node[box, right=of a] (b) {Xử lý};
  \node[box, right=of b] (c) {Đầu ra};
  \draw[-{Latex[length=2mm]}] (a) -- (b);
  \draw[-{Latex[length=2mm]}] (b) -- (c);
\end{tikzpicture}
\end{center}
\end{frame}

%---------------------------------------------------------
\section{Kết luận}
%---------------------------------------------------------

\begin{frame}
\frametitle{Kết luận}
\begin{itemize}
  \item Đóng góp 1
  \item Đóng góp 2
\end{itemize}
\vspace{0.5cm}
\begin{center}
{\Large \textbf{Cảm ơn quý thầy cô đã lắng nghe!}}
\end{center}
    \label{endmain}   % mốc hết phần chính; footer đếm tới đây
\end{frame}

%---------------------------------------------------------
% Slide dự phòng: đánh số riêng qua \appendix, không tính vào phần chính.
\appendix

\begin{frame}
\frametitle{Dự phòng --- câu hỏi thường gặp}
Nội dung trả lời câu hỏi.
\end{frame}

\end{document}
